\chapter{References}
\label{app:references}

\renewcommand{\bibname}{}
\nocite{*}

% Drop the underline from \url here only: underlined text cannot be broken
% across lines, which made the longer URLs below overflow the right margin.
% The thesis body keeps the underlined style.
\makeatletter
\renewcommand{\url}[1]{\ltx@url{#1}}

% book.cls opens thebibliography with its own \chapter*, and every \chapter in
% this class starts with \cleardoublepage. On top of the \chapter{References}
% above that meant two chapter openings in a row: the second one had to land on
% an odd page, leaving a blank one between the title and the list. Swallow the
% starred \chapter (and the running head it would reset, so these pages keep the
% chapter's own) -- nothing follows the bibliography, so this stays local to it.
\let\bib@chapter\chapter
\def\chapter{\@ifstar{\@gobble}{\bib@chapter}}
\let\@mkboth\@gobbletwo
\makeatother

\bibliographystyle{plain}
\bibliography{references}

\section{Additional References for the Glossary}
% This list is only for sources that are NOT already in the main References
% above (use plain \cite{} in glossary-entries.tex for those -- sharing a key
% between the two lists makes LaTeX print the wrong citation number, since
% \citegloss{} funnels into the same global, key-keyed "already cited"
% tracking as \cite{}). \citegloss{} is for glossary-only sources with no
% citation anywhere else in the thesis. \chapter/\@mkboth are already
% patched above; this being the true last chapter, no re-patching needed.
\bibliographystylegloss{plain}
\glosscontinuenumtrue
\bibliographygloss{references-glossary}
\glosscontinuenumfalse
